\section{F-SM-33: a countable-coproduct refinement}

\subsection{Statement}
Keep the abelian category, enough projectives, scalar adjunction,
exactness and coproduct preservation of \(\kappa\), and full
acyclic-monomorphism assumption from F-SM-32.  Replace its
coproduct and scalar-normalization assumptions by:
\begin{enumerate}
\item for every countable family \((X_n)\), the canonical map
\[
       \bigoplus_n F(X_n)\longrightarrow F(\bigoplus_nX_n)
                         \quad\text{is an equivalence};       \tag{SM33.1}
\]
\item for every small set \(I\), the canonical derived counit
\[
                  F(\kappa k^{(I)})\longrightarrow k^{(I)}[0]
                         \quad\text{is an equivalence}.       \tag{SM33.2}
\]
\end{enumerate}
Then \(F\) again induces
\[
                   \mathcal M[W^{-1}]_\infty\simeq D_{\ge0}(k)_\infty.
                                                               \tag{SM33.3}
\]
Exact countable coproducts in \(\mathcal M\) suffice for
(SM33.1), by the projective-resolution argument in
F-SM-32, Step~1.  No assertion that the scalar adjoints imply
(SM33.2) is included.

\subsection{Complete proof}
The finite cofibrant-object structure of F-SM-32, Step~2,
uses only short exact sequences and the factorization
\(X\hookrightarrow A_X\oplus Y\twoheadrightarrow Y\).
It therefore remains valid here.  By (SM33.1), \(W\) is
closed under countable coproducts.  Apply
Ramzi--Sosnilo--Winges Theorem~2.6(4) with the cardinal bound
\(\aleph_1\): the localization \(\mathcal C\) has countable
coproducts and its localization functor \(\gamma\) preserves
them.  It also has finite colimits.  The actual coherent
\(\overline F:\mathcal C\to\mathcal D_k\) preserves finite
colimits by the exact-triangle argument in F-SM-32, and
preserves countable coproducts by (SM33.1).

HTT Propositions~4.4.2.6--7, with the same cardinal bound,
give countable colimits and their preservation.  The nerve
of \(\Delta^{\mathrm{op}}\) is countable: the simplex
category has countably many objects and morphisms, and its
simplices are finite composable strings.  Hence
\(\mathcal C\) admits geometric realizations and
\(\overline F\) preserves them.  No uncountable colimit in
\(\mathcal C\) is needed below.

The left-fraction argument in F-SM-32, equation (SM32.6),
does not use any coproduct hypothesis and proves that
\(\overline F\) is conservative.  Naturality of the
canonical counit in (SM33.2) passes its equivalence from
free scalar modules to their projective retracts.
Exactness of \(\kappa\) again makes \(\varphi P\)
projective for projective \(P\in\mathcal M\).
Consequently the triangle identity makes
\(P\to\kappa\varphi P\) a member of \(W\).

Apply HA Theorem~1.3.3.8 to extend
\(Q\mapsto\gamma\kappa Q\) on projective scalars to a
realization-preserving functor \(L:\mathcal D_k\to\mathcal C\),
and \(P\mapsto\gamma P\) on projectives of \(\mathcal M\)
to \(T:\mathcal D_{\mathcal M}\to\mathcal C\).
All required targets admit geometric realizations by the
preceding paragraph.  The projective counits and units give,
by the full faithfulness of the same theorem,
\[
                 \overline F L\simeq\mathrm{id},\qquad
                         T\simeq L\mathbb L\varphi.           \tag{SM33.4}
\]
Here \(\mathbb L\varphi\) preserves realizations by
HA Lemma~1.3.3.11(2).

The completely proved comma-category lemma (SM32.12)--(SM32.13)
requires only geometric realizations in its two component
categories.  Apply it to
\((\mathrm{id}_{\mathcal C}\downarrow\gamma)\).
Extending the functor \(P\mapsto(\gamma P,P,\mathrm{id})\)
gives the natural transformation \(b:T\to\gamma H_0\).
After applying \(\overline F\), the arrow-valued uniqueness
in HA Theorem~1.3.3.8 identifies this with
\[
          \mathbb L\varphi(C)\longrightarrow
                          \mathbb L\varphi(H_0(C)[0]).
\]
This is precisely the comparison from F-SM-32, Step~5;
its proof uses only realization preservation and is
applicable with the countable colimits just established.
On \(C=X[0]\), the last arrow is the identity truncation.
Conservativity therefore gives \(T\iota\simeq\gamma\).
Combine this with (SM33.4) and descend the resulting natural
equivalence \(L\overline F\gamma\simeq\gamma\) through
the relative localization.  It yields
\(L\overline F\simeq\mathrm{id}_{\mathcal C}\), proving
(SM33.3). \(\square\)

\subsection{F-SM-35: when an exact tensor convention already forces AB4}
Let \(\mathcal A\) be an abelian category with small products
and coproducts.  Assume that for every small set \(I\), the
copower functor \(X\mapsto X^{(I)}\) preserves monomorphisms.
Then \(\mathcal A\) is AB4.

\paragraph{Proof.}
Let \((f_i:A_i\hookrightarrow B_i)_{i\in I}\) be any
family of monomorphisms.  The map
\(f=\prod_i f_i:\prod_iA_i\to\prod_iB_i\) is monic,
because products preserve kernels.  Therefore
\(f^{(I)}\) is monic by hypothesis.  Put
\(A=\prod_iA_i\), \(B=\prod_iB_i\).  The coordinate
inclusions \(e_i:A_i\to A\), with every other coordinate
zero, and projections \(p_i:A\to A_i\) define maps
\[
 \bigoplus_iA_i\xrightarrow{e}A^{(I)}
                   \xrightarrow{p}\bigoplus_iA_i,
 \qquad pe=1.
\]
Here \(e\) sends the \(i\)-th summand into the \(i\)-th
copy of \(A\), and \(p\) sends that copy through \(p_i\)
into the \(i\)-th summand.  The identical construction for
\(B\) commutes with \(\bigoplus_i f_i\) and \(f^{(I)}\).
Thus \(\bigoplus_i f_i\) is a retract of the monomorphism
\(f^{(I)}\) in the arrow category, so is monic.
Coproducts already preserve cokernels, being left adjoints
to diagonal functors.  They are therefore exact. \(\square\)

In particular, if \(\mathcal A\) has an exact tensor and
natural identifications
\[
              X\otimes\kappa(k^{(I)})\simeq X^{(I)},
                                                               \tag{SM35.1}
\]
the lemma applies.  A closed symmetric monoidal structure
supplies (SM35.1), because tensoring preserves coproducts
and \(\kappa(k^{(I)})=\mathbf1^{(I)}\).
This is a conditional interpretation result: the source
pages inspected say bilinear and exhibit the scalar
adjoints; they do not explicitly supply this full closed,
exact tensor convention or all products.  We consequently
do not use the lemma to declare bare (a)--(d) proved.

\subsection{Primary references}
The known inputs, with their exact primary locators and URLs,
are those listed in F-SM-32: Ramzi, Sosnilo, and Winges
(2025), Definition~2.5 and Theorem~2.6, pp.~5--6,
\url{https://arxiv.org/abs/2503.11338};
Cisinski (2019, author version July~30, 2026),
Proposition~7.5.6, pp.~372--373,
\url{https://cisinski.app.uni-regensburg.de/CatLR.pdf};
Lurie (2017), Theorem~1.3.3.8, p.~101, and
Lemma~1.3.3.11(2), pp.~103--104,
\url{https://people.math.harvard.edu/~lurie/papers/HA.pdf};
and Lurie (2009), Propositions~4.4.2.6--7, pp.~296--297,
\url{https://www.math.ias.edu/~lurie/papers/HTT.pdf}.
F-SM-35 is proved directly above.
